% =============================================================
% latex/ean13-encoding.tex
%
% Lookup-Tabellen + parametrisierter EAN-13-Renderer.
% Stellt die Befehle bereit:
%
%   \eanbarcode{<13-stelliger EAN>}
%       rendert den Strichcode als TikZ-Bild (Modulbreite 1mm,
%       Strichhöhe 30 Module für Daten und 33 für Guards,
%       Klarschrift in Source Sans 3).
%
%   \eanbits{<13-stelliger EAN>}{\MACRO}
%       schreibt die komma-separierte 95-Bit-Liste in \MACRO,
%       für Wiederverwendung in der Zeichenvorlage.
%
%   \eantableset{<EAN>}{<pos>}   gibt 'L', 'G' oder 'R' an Position <pos>.
%                                Position 1 ist die Erstziffer (kein Strich),
%                                ergibt '-'. Position 2..7 ist die linke
%                                Hälfte (L oder G), Position 8..13 immer R.
%   \eantablebits{<EAN>}{<pos>}  gibt die 7 Bits an Position <pos>
%                                als String '0001101' (Position 1: leer).
%
% Voraussetzungen: TikZ und xstring (in praeambel.tex geladen).
% =============================================================

\makeatletter

% --- L/G/R-Codes als 7-Bit-Komma-Listen ---
\expandafter\def\csname eanL@0\endcsname{0,0,0,1,1,0,1}
\expandafter\def\csname eanL@1\endcsname{0,0,1,1,0,0,1}
\expandafter\def\csname eanL@2\endcsname{0,0,1,0,0,1,1}
\expandafter\def\csname eanL@3\endcsname{0,1,1,1,1,0,1}
\expandafter\def\csname eanL@4\endcsname{0,1,0,0,0,1,1}
\expandafter\def\csname eanL@5\endcsname{0,1,1,0,0,0,1}
\expandafter\def\csname eanL@6\endcsname{0,1,0,1,1,1,1}
\expandafter\def\csname eanL@7\endcsname{0,1,1,1,0,1,1}
\expandafter\def\csname eanL@8\endcsname{0,1,1,0,1,1,1}
\expandafter\def\csname eanL@9\endcsname{0,0,0,1,0,1,1}

\expandafter\def\csname eanG@0\endcsname{0,1,0,0,1,1,1}
\expandafter\def\csname eanG@1\endcsname{0,1,1,0,0,1,1}
\expandafter\def\csname eanG@2\endcsname{0,0,1,1,0,1,1}
\expandafter\def\csname eanG@3\endcsname{0,1,0,0,0,0,1}
\expandafter\def\csname eanG@4\endcsname{0,0,1,1,1,0,1}
\expandafter\def\csname eanG@5\endcsname{0,1,1,1,0,0,1}
\expandafter\def\csname eanG@6\endcsname{0,0,0,0,1,0,1}
\expandafter\def\csname eanG@7\endcsname{0,0,1,0,0,0,1}
\expandafter\def\csname eanG@8\endcsname{0,0,0,1,0,0,1}
\expandafter\def\csname eanG@9\endcsname{0,0,1,0,1,1,1}

\expandafter\def\csname eanR@0\endcsname{1,1,1,0,0,1,0}
\expandafter\def\csname eanR@1\endcsname{1,1,0,0,1,1,0}
\expandafter\def\csname eanR@2\endcsname{1,1,0,1,1,0,0}
\expandafter\def\csname eanR@3\endcsname{1,0,0,0,0,1,0}
\expandafter\def\csname eanR@4\endcsname{1,0,1,1,1,0,0}
\expandafter\def\csname eanR@5\endcsname{1,0,0,1,1,1,0}
\expandafter\def\csname eanR@6\endcsname{1,0,1,0,0,0,0}
\expandafter\def\csname eanR@7\endcsname{1,0,0,0,1,0,0}
\expandafter\def\csname eanR@8\endcsname{1,0,0,1,0,0,0}
\expandafter\def\csname eanR@9\endcsname{1,1,1,0,1,0,0}

% --- Paritätsmuster der Erstziffer (als String LLGLGG etc.) ---
\expandafter\def\csname eanP@0\endcsname{LLLLLL}
\expandafter\def\csname eanP@1\endcsname{LLGLGG}
\expandafter\def\csname eanP@2\endcsname{LLGGLG}
\expandafter\def\csname eanP@3\endcsname{LLGGGL}
\expandafter\def\csname eanP@4\endcsname{LGLLGG}
\expandafter\def\csname eanP@5\endcsname{LGGLLG}
\expandafter\def\csname eanP@6\endcsname{LGGGLL}
\expandafter\def\csname eanP@7\endcsname{LGLGLG}
\expandafter\def\csname eanP@8\endcsname{LGLGGL}
\expandafter\def\csname eanP@9\endcsname{LGGLGL}

% --- Bit-Stream-Zusammenbau ---
% \eanbits{<13>}{\OUT}: schreibt die 95-Bit-Komma-Liste in \OUT.
\NewDocumentCommand{\eanbits}{mm}{%
  \StrChar{#1}{1}[\@eanFirstD]%
  \edef\@eanPattern{\csname eanP@\@eanFirstD\endcsname}%
  % Start-Guard
  \xdef#2{1,0,1}%
  % Linke 6 Ziffern, Pattern aus \@eanPattern auswählen
  \foreach \@p in {1,...,6}{%
    \StrChar{#1}{\the\numexpr\@p+1\relax}[\@eanZ]%
    \StrChar{\@eanPattern}{\@p}[\@eanS]%
    \xdef#2{#2,\csname ean\@eanS @\@eanZ\endcsname}%
  }%
  % Mittel-Guard
  \xdef#2{#2,0,1,0,1,0}%
  % Rechte 6 Ziffern (immer R)
  \foreach \@p in {1,...,6}{%
    \StrChar{#1}{\the\numexpr\@p+7\relax}[\@eanZ]%
    \xdef#2{#2,\csname eanR@\@eanZ\endcsname}%
  }%
  % End-Guard
  \xdef#2{#2,1,0,1}%
}

% --- Renderer ---
% \eanbarcodebits{<bits>}{<first>}{<left6>}{<right6>}: rendert eine
% vorberechnete Bit-Liste. Daten-Striche y=0..33, Guard-Striche y=-3..33.
\newcommand{\eanbarcodebits}[4]{%
  \begin{tikzpicture}[x=1mm, y=1mm, baseline=(current bounding box.south)]
    \foreach \bit [count=\j from 0] in {#1}{%
      \ifnum\bit=1
        \def\bo{0}%
        \ifnum\j<3\relax\def\bo{-3}\fi
        \ifnum\j>91\relax\def\bo{-3}\fi
        \ifnum\j>44\relax\ifnum\j<50\relax\def\bo{-3}\fi\fi
        \fill[black] (\j,\bo) rectangle ({\j+1},33);%
      \fi
    }%
    \node[anchor=north, font=\sffamily, inner sep=2pt] at (-5,-0.5) {#2};%
    \node[anchor=north, font=\sffamily, inner sep=2pt] at (24,-0.5) {#3};%
    \node[anchor=north, font=\sffamily, inner sep=2pt] at (71,-0.5) {#4};%
  \end{tikzpicture}%
}

% \eanbarcode{<13>}: bequeme Form, baut Bit-Liste und rendert.
\NewDocumentCommand{\eanbarcode}{m}{%
  \eanbits{#1}{\@eanFullBits}%
  \StrLeft{#1}{1}[\@eanFirst]%
  \StrMid{#1}{2}{7}[\@eanLeft]%
  \StrMid{#1}{8}{13}[\@eanRight]%
  \expandafter\eanbarcodebits\expandafter{\@eanFullBits}{\@eanFirst}{\@eanLeft}{\@eanRight}%
}

% --- Hilfen für die Zeichenvorlagen-Tabelle ---
% \eantableset{<EAN>}{<pos>} → 'L', 'G', 'R' oder '-'
\NewDocumentCommand{\eantableset}{mm}{%
  \ifnum#2=1 -\else
    \ifnum#2>7 R\else
      \StrChar{#1}{1}[\@eanFirstD]%
      \edef\@eanPattern{\csname eanP@\@eanFirstD\endcsname}%
      \StrChar{\@eanPattern}{\the\numexpr#2-1\relax}[\@eanS]%
      \@eanS%
    \fi
  \fi
}

% \eantablebits{<EAN>}{<pos>} → 7-Bit-String '0001101' oder leer
\NewDocumentCommand{\eantablebits}{mm}{%
  \ifnum#2=1 \else
    \ifnum#2>7
      \StrChar{#1}{#2}[\@eanZ]%
      \@eanCSToStr{eanR@\@eanZ}%
    \else
      \StrChar{#1}{1}[\@eanFirstD]%
      \edef\@eanPattern{\csname eanP@\@eanFirstD\endcsname}%
      \StrChar{\@eanPattern}{\the\numexpr#2-1\relax}[\@eanS]%
      \StrChar{#1}{#2}[\@eanZ]%
      \@eanCSToStr{ean\@eanS @\@eanZ}%
    \fi
  \fi
}

% Hilfsmacro: csname-Inhalt (Komma-Liste) als String ohne Kommata ausgeben
\newcommand{\@eanCSToStr}[1]{%
  \edef\@tmpcs{\csname #1\endcsname}%
  \StrSubstitute{\@tmpcs}{,}{}%
}

\makeatother
